\section{Introduction}\label{sec:intro}
A generative model learns, from samples of an unknown distribution $p_1$ on $\mathbb R^d$, to draw new samples from it. Diffusion models \citep{ho2020ddpm,song2021sde} and flow matching \citep{lipman2023fm,liu2023rf,albergo2023si} do this by transporting Gaussian noise $x_0\sim\mathcal N(0,I)$ to data $x_1\sim p_1$. Conditional flow matching (CFM) makes training a least-squares regression: draw a pair $(x_0,x_1)$ and a time $t\in[0,1]$, and regress a network $v_\theta(x_t,t)$ on the velocity $x_1-x_0$ of the straight line through the pair. Sampling integrates the ODE $\dot x=v_\theta(x,t)$ from $t=0$ to $1$, at a cost of one network evaluation (NFE) per Euler step. The learned paths are not straight, so few-step samples are poor: one Euler step of the model trained on independent pairs (I-CFM) maps every noise vector to almost the same blurry image (Figure~\ref{fig:samples}). Obtaining good samples with few NFE is therefore the central practical question, as it was for diffusion, where it led from DDPM's 1,000 steps to the samplers of \citet{song2021ddim} and \citet{karras2022edm}.

In our class presentation we derived CFM from the continuity equation and showed on four 2D datasets that pairing noise and data by minibatch optimal transport \citep[OT-CFM;][]{tong2024otcfm} makes the paths nearly straight, so that one Euler step is already close to the data ($W_2=0.12$--0.21, against 1.28--1.37 for I-CFM). On MNIST the advantage shrank: OT-CFM was better only at one to five steps, and in 784 dimensions minibatch OT saved 4.5\% of the pairing cost, against 87--93\% in 2D. This report asks why, and what reduces the NFE instead. Since diffusion with Gaussian noise and independent pairs is the same regression as I-CFM in other coordinates (Section~\ref{sec:diffusion}), the sampling problem is shared by both families, while the pairing is something only CFM can change.

With one network and one training budget fixed, we ask: \textbf{RQ1}, how reliably can we tell that one model or sampler is better, and what does FID miss? \textbf{RQ2}, at a fixed NFE, how much do the solver and the time grid matter? \textbf{RQ3}, how much do training-side changes (larger OT batches, the training-time distribution, reflow) reduce the NFE, and at what cost? \textbf{RQ4}, do the two sides combine, and does path straightness predict few-step quality? Each question is analysed theoretically first (Section~\ref{sec:theory}) and then measured against an explicit estimate of the evaluation error in six experiments, C1--C6 (Sections~\ref{sec:design} and~\ref{sec:experiments}), using about 20 GPU-hours of training.

\paragraph{Contributions.}
(1) A controlled comparison of training- and inference-side ways to reduce the NFE of CFM with one network and one budget, reporting the training cost of each and judging every gain against the seed-to-seed spread of FID.
(2) A time grid with step density $\propto\sqrt\kappa$, derived from the local Euler error with a gain predicted before sampling, and the finding that second-order solvers fail on CFM models because they evaluate the network at $t=1$, where its ideal value is known in closed form but badly learned; changing only the last step fixes this and lowers the NFE needed for FID 20 by 28--58\%.
(3) A bound, $2\operatorname{tr}\Sigma^{1/2}$, on what any pairing can save, explaining why minibatch OT helps little in high dimensions, and evidence that the training-time distribution matters more than the pairing.
(4) A quantification of evaluation error and a class-coverage test ($k$-means with a two-sample $\chi^2$ test) that separates image quality from class balance and shows where FID alone misleads.
